% =============================================================================
% Chapter: Privacy - K-Anonymity & L-Diversity
% =============================================================================

\chapter{Privacy: K-Anonymity \& L-Diversity}

\section{Εισαγωγή}

\keyword{Data Privacy} αφορά την προστασία προσωπικών δεδομένων κατά τη δημοσίευση ή ανάλυση datasets. Ακόμα κι αν αφαιρεθούν τα ονόματα, άλλα attributes μπορούν να οδηγήσουν σε \keyword{re-identification}.

\begin{warningbox}[Το Πρόβλημα]
\begin{center}
\begin{tikzpicture}
    % Published dataset
    \node[draw=primary, fill=box-blue, rounded corners=5pt, minimum width=4.5cm, minimum height=2.5cm] (pub) at (0,0) {};
    \node[anchor=north, font=\bfseries] at (pub.north) {Published Data};
    \node[font=\small] at (pub.center) {
        \begin{tabular}{cc}
            Age & Disease \\
            28 & HIV \\
            28 & Cancer
        \end{tabular}
    };

    % External knowledge
    \node[draw=secondary, fill=secondary!20, rounded corners=5pt, minimum width=4.5cm, minimum height=2.5cm] (ext) at (6,0) {};
    \node[anchor=north, font=\bfseries] at (ext.north) {External Knowledge};
    \node[font=\small, align=center] at (ext.center) {
        "John is 28 years old\\
        and lives in zipcode 12345"
    };

    % Arrow and result
    \draw[->, >=Stealth, thick, accent] (pub.east) -- (ext.west) node[midway, above] {+};

    % Result
    \node[draw=accent, fill=accent!20, rounded corners=5pt, minimum width=4cm, minimum height=1.5cm] (res) at (3,-2.5) {};
    \node[align=center, font=\small] at (res.center) {
        \textbf{Re-identification!}\\
        Can identify John's disease
    };
    \draw[->, >=Stealth, thick, accent] (3,-0.8) -- (res.north);
\end{tikzpicture}
\end{center}
\end{warningbox}

% =============================================================================
\section{Attribute Types}

\begin{center}
\begin{tikzpicture}
    % Explicit Identifiers
    \node[draw=accent, fill=accent!20, rounded corners=10pt, minimum width=5cm, minimum height=2.5cm] (ei) at (0,3) {};
    \node[anchor=north, font=\bfseries\color{accent}] at (ei.north) {Explicit Identifiers};
    \node[align=center, font=\small] at (ei.center) {
        Μοναδικά αναγνωρίζουν\\
        \textbf{Name, SSN, Phone}\\
        $\Rightarrow$ ΑΦΑΙΡΟΥΝΤΑΙ πάντα
    };

    % Quasi-Identifiers
    \node[draw=warning, fill=warning!20, rounded corners=10pt, minimum width=5cm, minimum height=2.5cm] (qi) at (6.5,3) {};
    \node[anchor=north, font=\bfseries\color{warning}] at (qi.north) {Quasi-Identifiers (QIs)};
    \node[align=center, font=\small] at (qi.center) {
        Σε \textbf{συνδυασμό} ταυτοποιούν\\
        \textbf{Age, Zipcode, Gender}\\
        $\Rightarrow$ Γενικεύονται
    };

    % Sensitive Attributes
    \node[draw=secondary, fill=secondary!20, rounded corners=10pt, minimum width=5cm, minimum height=2.5cm] (sa) at (0,0) {};
    \node[anchor=north, font=\bfseries\color{secondary}] at (sa.north) {Sensitive Attributes};
    \node[align=center, font=\small] at (sa.center) {
        Πληροφορία προς προστασία\\
        \textbf{Disease, Salary}\\
        $\Rightarrow$ Προστατεύονται
    };

    % Non-Sensitive
    \node[draw=success, fill=success!20, rounded corners=10pt, minimum width=5cm, minimum height=2.5cm] (ns) at (6.5,0) {};
    \node[anchor=north, font=\bfseries\color{success}] at (ns.north) {Non-Sensitive};
    \node[align=center, font=\small] at (ns.center) {
        Δεν αποκαλύπτουν τίποτα\\
        \textbf{Favorite color}\\
        $\Rightarrow$ Δημοσιεύονται
    };
\end{tikzpicture}
\end{center}

% =============================================================================
\section{K-Anonymity}

\begin{formulabox}[K-Anonymity - Ορισμός]
Ένα dataset είναι \textbf{K-ANONYMOUS} αν για κάθε record, υπάρχουν τουλάχιστον $(k-1)$ άλλα records με \textbf{ΙΔΙΕΣ} τιμές στα Quasi-Identifiers.

\vspace{0.5em}
Με άλλα λόγια: κάθε \keyword{equivalence class} (συνδυασμός QIs) έχει $\geq k$ records.
\end{formulabox}

\begin{examplebox}[K-Anonymity Example]
\textbf{Original Data (NOT k-anonymous)}:
\begin{center}
\begin{tabular}{cc|c}
\toprule
Age & Zipcode & Disease \\
\midrule
25 & 12345 & HIV \\
27 & 12345 & Cancer \\
35 & 54321 & Flu \\
36 & 54321 & Diabetes \\
\bottomrule
\end{tabular}
\end{center}

\textbf{2-Anonymous (k=2)} - after generalization:
\begin{center}
\begin{tabular}{cc|c|l}
\toprule
Age & Zipcode & Disease & Class \\
\midrule
\rowcolor{success!20} 20-30 & 123** & HIV & \multirow{2}{*}{Class 1: 2 records \checkmark} \\
\rowcolor{success!20} 20-30 & 123** & Cancer & \\
\rowcolor{info!20} 30-40 & 543** & Flu & \multirow{2}{*}{Class 2: 2 records \checkmark} \\
\rowcolor{info!20} 30-40 & 543** & Diabetes & \\
\bottomrule
\end{tabular}
\end{center}
\end{examplebox}

\subsection{Anonymization Techniques}

\begin{center}
\begin{tikzpicture}
    % Generalization
    \node[draw=secondary, fill=secondary!20, rounded corners=10pt, minimum width=6cm, minimum height=3cm] (gen) at (0,0) {};
    \node[anchor=north, font=\bfseries\color{secondary}] at (gen.north) {Generalization};
    \node[align=center, font=\small] at (gen.center) {
        Make values more general\\[0.5em]
        Age: $25 \to 20\text{-}30 \to 20\text{-}40 \to *$\\
        Zip: $12345 \to 1234* \to 123** \to *****$
    };

    % Suppression
    \node[draw=accent, fill=accent!20, rounded corners=10pt, minimum width=6cm, minimum height=3cm] (sup) at (8,0) {};
    \node[anchor=north, font=\bfseries\color{accent}] at (sup.north) {Suppression};
    \node[align=center, font=\small] at (sup.center) {
        Remove values entirely\\[0.5em]
        Age: $25 \to *$\\
        Or remove entire records
    };
\end{tikzpicture}
\end{center}

\subsection{Αδυναμίες K-Anonymity}

\begin{warningbox}[Homogeneity Attack]
\begin{center}
\begin{tabular}{cc|c}
\toprule
Age & Zipcode & Disease \\
\midrule
\rowcolor{accent!20} 20-30 & 123** & HIV \\
\rowcolor{accent!20} 20-30 & 123** & HIV \\
\rowcolor{accent!20} 20-30 & 123** & HIV \\
\bottomrule
\end{tabular}
\end{center}

3-anonymous \textbf{αλλά}: αν ξέρω κάποιον 25χρονο στο 12345, \important{ΞΕΡΩ ότι έχει HIV!}

(Όλοι στο equivalence class έχουν την ίδια disease)
\end{warningbox}

% =============================================================================
\section{L-Diversity}

\begin{formulabox}[L-Diversity - Ορισμός]
Ένα dataset είναι \textbf{L-DIVERSE} αν σε κάθε equivalence class υπάρχουν τουλάχιστον $L$ \textbf{ΔΙΑΦΟΡΕΤΙΚΕΣ} τιμές του sensitive attribute.
\end{formulabox}

\begin{examplebox}[L-Diversity Example]
\textbf{2-Diverse}:
\begin{center}
\begin{tabular}{cc|c|l}
\toprule
Age & Zipcode & Disease & Distinct values \\
\midrule
\rowcolor{success!20} 20-30 & 123** & HIV & \multirow{2}{*}{\{HIV, Cancer\} = 2 \checkmark} \\
\rowcolor{success!20} 20-30 & 123** & Cancer & \\
\rowcolor{success!20} 30-40 & 543** & Flu & \multirow{2}{*}{\{Flu, Diabetes\} = 2 \checkmark} \\
\rowcolor{success!20} 30-40 & 543** & Diabetes & \\
\bottomrule
\end{tabular}
\end{center}

\textbf{NOT 2-Diverse}:
\begin{center}
\begin{tabular}{cc|c|l}
\toprule
Age & Zipcode & Disease & Distinct values \\
\midrule
\rowcolor{accent!20} 20-30 & 123** & HIV & \multirow{2}{*}{\{HIV\} = 1 \texttimes} \\
\rowcolor{accent!20} 20-30 & 123** & HIV & \\
\bottomrule
\end{tabular}
\end{center}
\end{examplebox}

% =============================================================================
\section{T-Closeness}

\begin{formulabox}[T-Closeness - Ορισμός]
\textbf{T-CLOSENESS}: Η κατανομή του sensitive attribute σε κάθε equivalence class πρέπει να είναι "κοντά" στη συνολική κατανομή (distance $\leq t$).

\textbf{Distance}: Earth Mover's Distance (EMD)
\end{formulabox}

\begin{warningbox}[Γιατί χρειάζεται - Skewness Attack]
\textbf{Overall distribution}: 99\% "Healthy", 1\% "Disease"

\textbf{Equivalence class}: 50\% "Healthy", 50\% "Disease"

L=2 satisfied, \textbf{BUT}: αν κάποιος είναι στο class, η πιθανότητα νόσου είναι \important{50\% αντί για 1\%!}
\end{warningbox}

% =============================================================================
\section{Σύγκριση Μέτρων}

\begin{center}
\begin{tikzpicture}
    \matrix (m) [
        matrix of nodes,
        nodes={
            draw=primary,
            minimum width=3.5cm,
            minimum height=1.2cm,
            anchor=center,
            font=\small
        },
        column 1/.style={nodes={fill=ntua-blue!20, font=\bfseries}},
        row 1/.style={nodes={fill=secondary!30, font=\bfseries}},
        column sep=-\pgflinewidth,
        row sep=-\pgflinewidth
    ] {
        Measure & Protects against & Weakness \\
        K-Anonymity & Re-identification & Homogeneity attack \\
        L-Diversity & Homogeneity & Skewness attack \\
        T-Closeness & Distribution leakage & Complex to implement \\
    };
\end{tikzpicture}
\end{center}

% =============================================================================
\section{Συχνά Λάθη}

\begin{warningbox}[Κοινά Λάθη στις Εξετάσεις]
\begin{enumerate}
    \item \textbf{K-Anonymity}:
    \begin{itemize}
        \item[$\times$] k-anonymity προστατεύει πλήρως
        \item[\checkmark] Προστατεύει μόνο από re-identification μέσω QIs
    \end{itemize}

    \item \textbf{QI Definition}:
    \begin{itemize}
        \item[$\times$] Age μόνο του είναι identifier
        \item[\checkmark] Age + Zipcode + Gender ΜΑΖΙ = quasi-identifier
    \end{itemize}

    \item \textbf{L-Diversity}:
    \begin{itemize}
        \item[$\times$] L=3 σημαίνει ακριβώς 3 values
        \item[\checkmark] L=3 σημαίνει ΤΟΥΛΑΧΙΣΤΟΝ 3 distinct values
    \end{itemize}
\end{enumerate}
\end{warningbox}

% =============================================================================
\section{Σύνοψη - Key Points}

\begin{center}
\begin{tikzpicture}
    \node[draw=ntua-blue, fill=ntua-blue!10, rounded corners=5pt, minimum width=14cm, minimum height=9cm] (main) at (0,0) {};

    % Attribute Types
    \node[anchor=north west, font=\bfseries\Large\color{ntua-blue}] at ([xshift=5pt, yshift=-5pt]main.north west) {Attribute Types};
    \node[anchor=north west, align=left, font=\small] at ([xshift=10pt, yshift=-30pt]main.north west) {
        \textbf{Explicit Identifiers} $\to$ REMOVE (Name, SSN)\\
        \textbf{Quasi-Identifiers} $\to$ GENERALIZE (Age, Zipcode)\\
        \textbf{Sensitive} $\to$ PROTECT (Disease, Salary)
    };

    % Privacy Measures
    \node[draw=secondary, fill=box-blue, rounded corners=5pt, minimum width=12cm, minimum height=4.5cm]
        at (0,-1.5) {};
    \node[anchor=north, font=\bfseries\color{secondary}] at (0,0.7) {Privacy Measures (ΑΠΟΜΝΗΜΟΝΕΥΣΕ!)};

    \node[align=left, font=\small] at (0,-1.5) {
        \textbf{K-ANONYMITY}:\\
        Every equivalence class has $\geq k$ records\\
        Weakness: Homogeneity attack\\[0.8em]
        \textbf{L-DIVERSITY}:\\
        Every equivalence class has $\geq L$ distinct sensitive values\\
        Weakness: Skewness attack\\[0.8em]
        \textbf{T-CLOSENESS}:\\
        Distribution in class $\approx$ overall distribution (distance $\leq t$)
    };

    % Techniques box
    \node[draw=success, fill=box-green, rounded corners=5pt, minimum width=6cm, minimum height=1.5cm]
        at (0,-4) {};
    \node[align=center, font=\small] at (0,-4) {
        \textbf{Techniques}\\
        Generalization: $25 \to 20\text{-}30$ | Suppression: $25 \to *$
    };
\end{tikzpicture}
\end{center}
